\chapter{Conclusion and Future Work}
\label{chap:conclusion}

\section{Answers to the Research Questions}
\label{sec:conclusion:1}

\paragraph{RQ1.} \emph{\rqOneText} \hybroid{} represents both sources at application level:
a 64-dimensional program-structure vector is concatenated with \netfeatures{} flow features
averaged per application. On CICAndMal2017, this fusion improved reported detection $F_1$ by one
or two percentage points over the better single source. For category
classification, the difference ranged from minus two to plus three points across learners.
These matched comparisons address gap~\gapref{1} within the evaluated corpus and protocol.

\paragraph{RQ2.} \emph{\rqTwoText} \hgann{} groups methods by call neighbourhood, semantic
similarity and security-sensitive slices, and learns membership weights by attention.
On Drebin and CICMalDroid, five-run mean accuracy exceeded the stronger pairwise baseline by
3.1 to 10.8 points for every hypergraph configuration under common inputs and fixed splits.
\hgann{} led the strongest non-attentive baseline in three of four settings by 1.2 to 4.1
points; its CICMalDroid binary accuracy was 97.6\% against 97.7\%. These complete-system
comparisons address gap~\gapref{2}; attention attribution requires a controlled ablation.

\paragraph{RQ3.} \emph{\rqThreeText} \nadics{} separates feature preparation from
interchangeable learning components to support researcher-guided comparison and selection of
candidate models for a dataset. On the session-disjoint S7-300 evaluation, a random forest
achieved 94.50\% attack $F_1$ at a record-level false-positive rate of 0.484\%.
The GAN-based extension with 24 PCA components achieved 91.89\% timestamp recall at a
false-positive rate of 5.08\% on later SWaT test data. These separate evaluations address
gap~\gapref{3} through a modular framework and observation-specific analysis within the tested
industrial settings.

\paragraph{RQ4.} \emph{\rqFourText} SAFAIR translated attacker assumptions into a contest
protocol and a benchmark with task-specific interfaces. The implementation and scoring analysis
address the design and measurement aspects of gap~\gapref{4}. Any constant classifier has zero
accuracy loss under label-preserving attacks. In Vicomtech's histopathology benchmark,
the activation detector had zero loss under PGD but an attacked accuracy 12.9 percentage points
below adversarial training. Clean and attacked accuracy therefore need to be reported together,
under an explicit prediction objective and valid perturbation constraints. The four contest
rows do not establish a general relation between clean performance and robustness.

\section{Summary of Contributions}
\label{sec:conclusion:2}

The Android studies form the main contribution. \hybroid{} combines static program-graph
representations with observed network flows and measures the value of fusion on the same
applications. The candidate designed the overall pipeline and traffic analysis and conducted
the evaluation; his contribution to graph processing was partial.
\hgann{} represents context within the application package. Chapter~\ref{chap:hgann}
formalises its construction and shows that different hyperedge groupings can produce different
attentive propagation even when their fixed propagation matrices coincide.

The industrial contribution is a modular framework for semi-automated model development on
network and process observations. It combines configurable preparation with comparison of
candidate learners, supported by the candidate's extractor and labelled corpus and a GAN-based
extension for SWaT.
SAFAIR contributes a benchmark design for adversarial robustness evaluation and reusable
software with task-specific interfaces. Chapter~\ref{chap:advml} analyses accuracy-loss scoring and the effect of binary
label reversal after an alarm. Partner defences provide applications of the framework.

% Structure v0.3 placed a Limitations section here. It was removed at the author's instruction of
% 12 September 2026; the limits of each study stay in the Limitations sections of Chapters 3 to 6.

\section{Future Work}
\label{sec:conclusion:3}

The first priority is a controlled \hgann{} ablation on both corpora. Comparing learned and
fixed weighting over identical fused hyperedges requires common preprocessing and
normalisation, with reported variation across runs. A comparison with topological
hyperedges alone would assess the contribution of construction.
A common Android cohort with time-ordered splits would permit comparison of static relational
encodings and the additional value of traffic, as proposed in Chapter~\ref{chap:discussion}.
Acquisition effort and computation cost should be measured separately. Adversarial evaluation
would then apply valid package transformations that preserve malicious functionality, under
stated attacker access and budgets. It should measure clean and attacked performance of the
complete pipeline, including fusion.

The dissertation connects contextual representation design with the evidence and evaluation
procedures needed to assess learning-based security decisions.
